% =====================================================
% Chapter 5: Circuit Design
% General design methodology is presented; actual component
% values, topology choice, and calculations are project-specific
% and are left as clearly marked placeholders pending the
% team's final design.
% =====================================================
\chapter{Circuit Design}
\label{ch:circuitdesign}

\section{Design Methodology}

A general design procedure for an AGC microphone preamplifier, consistent with the architecture of Chapter~\ref{ch:architecture}, proceeds as follows \cite{ref1,ref2,ref4,ref5,ref6}:

\begin{enumerate}
    \item \textbf{Define the input/output level range} the circuit must accommodate (minimum and maximum expected microphone input, and the desired output level range), from the target specifications of Chapter~\ref{ch:problem}.
    \item \textbf{Design the fixed-gain (or nominal-gain) preamplifier stage}, selecting an op-amp (or transistor stage) with adequate bandwidth, low noise, and sufficient open-loop gain, and choosing the non-inverting or inverting gain-setting resistors using Eq.~\ref{eq:noninverting_gain} (or the corresponding inverting-gain relation) for the nominal (unity-control) gain.
    \item \textbf{Select and size the variable-gain element} (e.g., JFET-in-feedback, dedicated VCA IC, or diode-based attenuator, as discussed in Chapter~\ref{ch:theory}), and determine the control-voltage-to-gain relationship for the chosen device from its datasheet or characterisation data.
    \item \textbf{Design the level/envelope detector}, typically a rectifier (diode or precision rectifier) followed by a low-pass filter, choosing the filter time constant to realise the desired attack time.
    \item \textbf{Design the control-voltage generation and release-time filtering}, including any threshold-setting reference and compression-ratio-shaping network.
    \item \textbf{Verify DC bias points and power budget} for all active devices, ensuring operation within safe, low-voltage limits (Chapter~\ref{ch:safety}).
    \item \textbf{Iterate using simulation} (Chapter~\ref{ch:simulation}) before committing to hardware construction.
\end{enumerate}

\section{Final Circuit Schematic}

% =====================================================
% PLACEHOLDER -- ACTUAL PROJECT SCHEMATIC NOT YET AVAILABLE.
%
% This file intentionally does NOT contain a circuit schematic.
% No topology, component values, op-amp/transistor choices, or
% node labelling should be inferred from this placeholder.
%
% The team must replace this file with the final AGC
% microphone preamplifier schematic (exported at high resolution
% from Multisim, or redrawn in TikZ/Circuitikz), showing supply
% rails, all component values, and labelled test points, once
% the actual circuit design is finalised. See Chapter 5
% (Circuit Design) for the surrounding discussion.
% =====================================================
\figureplaceholder{%
    \textbf{Final circuit schematic not yet available.}\\[4pt]
    \footnotesize Insert the finalised AGC microphone preamplifier schematic here (exported from Multisim at high resolution), with supply rails, all component values, and test points labelled.%
}{Final circuit schematic and designated test points (placeholder -- to be replaced by the team's actual design).}
\label{fig:agc-schematic}


%% TODO: Once the schematic above is replaced with the team's actual
%% design, add a short walkthrough here describing signal flow through
%% the real circuit, referencing labelled test points (see
%% tables/test_points.tex).

\section{Design Calculations and Component Selection}

Every important resistor, capacitor, bias point, gain, and time-constant calculation for the implemented circuit should be shown in this section, following the general procedure above, together with the standard (E12/E24) component value actually selected and a brief justification. Table~\ref{tab:component-values} provides the structure for recording this; it is currently a placeholder because the team's specific calculations are not available in the materials supplied for this report.

% =====================================================
% Table: Major component selection
% Structure only -- component references, functions, and
% calculated/selected values are project-specific.
% =====================================================
\begin{table}[H]
\centering
\caption{Major component selection.}
\label{tab:component-values}
\begin{tabularx}{\textwidth}{@{}lXccc@{}}
\toprule
\textbf{Ref.} & \textbf{Function} & \textbf{Calculated} & \textbf{Selected} & \textbf{Rating/model} \\
\midrule
--- & --- & --- & --- & --- \\
--- & --- & --- & --- & --- \\
--- & --- & --- & --- & --- \\
\bottomrule
\end{tabularx}
\end{table}

\noindent\placeholder{Component references, functions, calculated values, selected standard values, and ratings/models to be completed once the circuit design is finalised}


\placeholder{Insert the worked calculations here: e.g. preamplifier gain calculation from Eq.~\ref{eq:noninverting_gain}, envelope-detector RC time-constant calculation for the target attack/release time, VCA control-voltage range calculation, and DC bias-point calculations, each with the resulting standard component value and its justification}

\section{Test Points}

Table~\ref{tab:test-points} lists the test points that should be defined on the schematic once finalised, to support the simulation (Chapter~\ref{ch:simulation}) and hardware verification (Chapters~\ref{ch:hardware}--\ref{ch:results}) stages.

% =====================================================
% Table: Test points
% Structure only -- actual test-point labels and expected
% signals depend on the team's final schematic.
% =====================================================
\begin{table}[H]
\centering
\caption{Designated test points.}
\label{tab:test-points}
\begin{tabularx}{\textwidth}{@{}lXX@{}}
\toprule
\textbf{Test point} & \textbf{Location / stage} & \textbf{Expected signal} \\
\midrule
--- & --- & --- \\
--- & --- & --- \\
--- & --- & --- \\
\bottomrule
\end{tabularx}
\end{table}

\noindent\placeholder{Test-point labels, locations, and expected signals to be defined on the finalised schematic}

